\documentclass[a4paper,12pt]{extarticle}
\usepackage[table,xcdraw]{xcolor}
\usepackage{pgfplots}
\usepackage[unicode, pdftex]{hyperref}
\usepackage[T2A]{fontenc}
\usepackage[utf8]{inputenc}
\usepackage[english,russian]{babel}
\usepackage{amsmath,amsfonts,amsthm, mathtools}
\usepackage{amssymb}
\usepackage{dsfont}
\usepackage{graphicx}
\graphicspath{{images/}}
\usepackage{wrapfig}
\usepackage{floatrow}
\usepackage{setspace}
\usepackage[headings]{fancyhdr}
\usepackage[margin=10pt,font={small,stretch=0,9},labelfont=bf,labelsep=period,
justification=centerlast]{caption}
\usepackage{array,tabularx,tabulary,booktabs}
\RequirePackage{multirow}
\RequirePackage{multicol}
\RequirePackage{longtable}
\usepackage{indentfirst}
\usepackage{hyperref}
\usepackage{icomma}
\newcommand*{\hm}[1]{#1\nobreak\discretionary{}
	{\hbox{$\mathsurround=0pt #1$}}{}}
\usepackage{soulutf8} 
\usepackage{etoolbox} 
\usepackage{geometry}
\geometry{top=20mm}
\geometry{bottom=20mm}
\geometry{left=20mm}
\geometry{right=20mm}
\DeclareMathOperator{\sgn}{\mathop{sgn}}
\usepackage[OT1]{fontenc}
\usepackage{amsmath}
\usepackage{amsfonts}
\usepackage{amssymb}
\usepackage{wasysym}
\usepackage{wrapfig}
\usepackage{physics}
\usepackage[normalem]{ulem}
\graphicspath{{pictures/}}
\DeclareGraphicsExtensions{.png,.jpg,.jpeg}
\usepackage{indentfirst}
\usepackage[T2A]{fontenc}
\usepackage{subfigure}
\usepackage{fancyhdr}
\usepackage{caption}

\usepackage{tabularx}
\usepackage{floatrow}
\usepackage{multicol}
\usepackage{lastpage}

\pgfplotsset{
  MyStyle1/.style={
    axis lines = middle,
    enlargelimits = true,
    grid=both,
    grid style={line width=.1pt, draw=gray!10},
    major grid style={line width=.2pt,draw=gray!50},
    minor tick num=5,
    enlargelimits=0.05,
    ticklabel style={font=\tiny,fill=white},
    xlabel style={at={(ticklabel* cs:1)},anchor=north west},
    ylabel style={at={(ticklabel* cs:1)},anchor=south east},
  }
}
\def\Slope{1.5}
\def\Intercept{-20}


% Настройка отчёркиваний и нумерации
\usepackage{fancyhdr}
\pagestyle{fancy}
\fancyhf{}

\fancyhead[L]{\nouppercase{\leftmark\hfill}}
\fancyfoot[RO, RE]{Страница \thepage \ из \pageref{LastPage}}
\renewcommand{\headrulewidth}{0.5pt}
\renewcommand{\footrulewidth}{0.5pt}

\usepackage{titlesec}
\titlelabel{\thetitle.\quad}

\pgfplotsset{compat=1.18}
\begin{document}

\thispagestyle{empty}
\begin{center}
	\textit{Федеральное государственное автономное образовательное\\ учреждение высшего образования }
	
	\vspace{0.5ex}
	
	\textbf{«Московский физико-технический институт\\ (национальный исследовательский университет)»}
\end{center}

\vspace{10ex}
\begin{figure}[!h]
  \begin{center}
    
\includegraphics[width = 0.8\textwidth]{MIPT.png}
  \end{center}
\end{figure}

\begin{center}
	
	\textbf{Лабораторная работа № 5.8.1}
	
	\vspace{1ex}
	
	по курсу общей физики
	
	на тему:
	
	\textbf{\textit{<<Тепловое излучение>>}}
	
	\vspace{20ex}
	
	\begin{flushright}
		\noindent
		\textit{Работу выполнил:}\\  
		\textit{Легких Никита \\(группа Б05-301)}
	\end{flushright}
	\vfill
	Долгопрудный \\ \today
	\setcounter{page}{1}
\end{center}
\newpage

% Конец шаблона
\newpage

\textbf{Цель работы:} При помощи модели абсолютно чёрного тела (АТЧ) провести измерения температуры оптическим пирометром с исчезающей нитью и термопарой, исследовать излучение накалённых тел с различной испускательной способностью, определить постоянные Планка и Стефана-Больцмана.

\textbf{Оборудование:} оптический пирометр (модель: Проминь М1), модель абсолютно чёрного тела, образцы колец, вольфрамовая лампа, неоновая лампа, блок питания, цифровые мультиметры (модель: GW Instek GDM-8145).

\section{Теоретические положения}
Для измерения температуры разогретых тел, удалённых от наблюдателя, применяют методы оптической пирометрии, основанные на использовании зависимости испускательной способности исследуемого тела от температуры. Различают три температуры, функционально связанные с истинной термодинамической температурой и излучательной способностью тела: радиационную $T_{rad}$, цветовую $T_{col}$ и яркостную $T_b_r$. \par
В работе измеряется яркостная температура. \textbf{Яркостная температура} - это температура абсолютно чёрного тела, при которой его спектральная испускательная способность равна спектральной испускательной способности исследуемого тела при той же длине волны.
 Измерение яркостной температуры раскалённого тела производится при помощи оптического пирометра с исчезающей нитью, основанного на визуальном сравнении яркости раскалённой нити с яркостью изображения исследуемого тела. \par
Яркостная температура тела всегда ниже его термодинамической температуры. Это связано с тем, что любое нечёрное тело излучает меньше, чем АЧТ при той же температуре. Зависимость между яркостной и термодинамической температурами вольфрама приведена на рис. 1

\begin{figure}[h]
    \centering
    \includegraphics[width=10cm]{pics/2.png}
    \caption{График зависимости $T = f(T_{br})$ для вольфрам}
    \label{fig:vac}
\end{figure}

По результатам измерений мощности излучения вольфрамовой нити можно судить о справедливости закона Стефана-Больцмана. Если бы нить излучала как АЧТ, то баланс потребляемой и излучаемой энергии определялся бы соотношением 
\begin{equation}
    W = \sigma S (T^4 - T_0^4),
\end{equation}
где $W$ - потребляемая нитью электрическая мощность, $S$ - площадь излучающей поверхности нити, $T$ - температура нити, $T_0$ - температура окружающей среды. Однако вольфрамовая нить излучает как серое тел, и излучение её ослаблено по сравнению с АЧТ в $\varepsilon_T$ раз для любой волны при данной температуре тела Т. Тогда предположив, что нить излучает как серое тело и с учётом того, что $T_0 \ll T$, выражение (1) можно переписать в виде
\begin{equation}
    W = \varepsilon_T S \sigma T^4
\end{equation}
В справедливости закона Стефана-Больцмана можно убедиться, построив график зависимости $W(T)$ в логарифмическом масштабе и по углу наклона определить показатель степени $n$ исследуемой температурной зависимости. В пределах погрешности показатель степени должен быть близок к четырём. \par
Также из формулы (2) можно определить постоянную Стефана-Больцмана.

\section*{Экспериментальная установка}
\begin{figure}[H]
    \centering
    \includegraphics[width=9.77cm]{pics/ustanovka.png}
    \caption{Схема экспериментальной установки: 1 - блок питания; 2 - тумблер включения питания образцов; 3 - тумблер нагрева нити пирометра; 4 - кнопка "Нагрев нити"; 5 - кнопка "охлаждение нити"; 6 - тумблер переключения образцов; 7 - регулятор мощности нагрева образцов; 8 - окуляр пирометра; 9 - корпус пирометра; 10 - объектив пирометра; 11 - переключение диапазонов; 12 - ручка смещения красного светофильтра; 13 - регулировочный винт; 14 - вольтметр (напряжение на лампе накаливания); 15 - амперметр (ток через образцы); 16 - вольтметр в цепи термопары; 17 - модель АЧТ; 18 трубка с кольцами из материалов с различной излучательной способностью; 19 - лампа накаливания; 20 - неоновая лампочка.}
\end{figure}
Для используемых в работе цифровых мультиметров погрешности считаются следующим образом в используемых нами диапазонах:
\begin{enumerate}
    \item 2V: $\sigma_V = \pm \left( 0.001 \cdot |V| + 0.002 \right)$
    \item 2A: $\sigma_V = \pm \left( 0.005 \cdot |I| + 0.002 \right)$
    \item 20V: $\sigma_V = \pm \left( 0.01 \cdot |V| + 0.05 \right)$
\end{enumerate}
\section*{Ход работы}
\subsection*{Изучение работы оптического пирометра}
С помощью пирометра измерим температуру модели АЧТ (повторим измерение трижды). Результаты внесём в таблицу ниже:
\begin{table}[H]
    \centering
    \begin{tabular}{|c|c|c|c|}
        \hline
        $T_1$, $\text{С}^{\circ}$ & $T_2$, $\text{С}^{\circ}$ & $T_3$, $\text{С}^{\circ}$ & $T_{\text{пир}}$, $\text{С}^{\circ}$ \\ \hline
        $1085 \pm 10$ & $1100 \pm 10$ & $1095 \pm 10$ & $1093 \pm 6$ \\ \hline
    \end{tabular}
    \caption{Результаты измерений пирометра и их среднее}
\end{table}
Также найдем значение температуры АЧТ с помощью хромель-алюминиевой термопары и цифрового вольтметра по формуле $T_{\text{тер}} = V/k + T_{\text{ком}}$, где $k = 41$ $\text{мкВ}/\text{С}^{\circ}$ - постоянная термопары. Получим:
\begin{center}
    $V = (43,8 \pm 0,5)$ мВ, $T_{\text{ком}} = (22 \pm 0,5)$ $\text{С}^{\circ}$ $\Rightarrow$ $T_{\text{тер}} = (1092\pm 13)$ $\text{С}^{\circ}$
\end{center}
Таким образом, значения температуры АЧТ, измеренные оптическим пирометром и термопарным термометром, совпадают в пределах доверительных интервалов, из чего можно сделать вывод о том, что измерения пирометра достаточно достоверны.
\subsection*{Измерение яркостной температуры накаленных тел}
Покажем, что разные тела, накалённые до одинаковой термодинамической температуры, имеют различную яркостную температуру. Для этого прогреем керамическую трубку с образцами до красного каления. Пронаблюдаем различную яркость поверхности трубки и колец:
\begin{figure}[H]
    \centering
    \includegraphics[width=5cm]{pics/1.jpg}
\end{figure}
Несовпадение яркостной температуры у различных тел, имеющих одинаковую термодинамическую температуру, вытекает из того, что эти две величины связаны, в том числе, через спектральный коэффициент поглощения, который у разных материалов различен.
\subsection*{Проверка закона Стефана-Больцмана}
Постепенно увеличивая накал нити лампы, начиная со слабого тёмно-красного накала до 1900$^{\circ}$C, будем измерять пирометром яркостную температуру нити, а также значение силы тока и напряжения на ней. Результаты измерений занесём в таблицу ниже:
\begin{table}[H]
\centering
\begin{tabular}{|c|c|c|c|c|}
\hline
$T_{\text{ярк}}$, $\text{С}^{\circ}$ & $I$, А & $\sigma_I$, А & $V$, В & $\sigma_V$, В \\ \hline
900        & 0,446& 0,005 & 1,44& 0,0022 \\ \hline
1000       & 0,494& 0,005 & 1,83& 0,0023 \\ \hline
1100       & 0,522& 0,005 & 2,09& 0,0024 \\ \hline
1200       & 0,561& 0,006 & 2,45& 0,0025 \\ \hline
1300       & 0,608& 0,006 & 2,89& 0,0025 \\ \hline
1400       & 0,667& 0,006 & 3,49& 0,0026 \\ \hline
1500       & 0,701& 0,007 & 3,86& 0,0027 \\ \hline
1600       & 0,781& 0,007 & 4,76& 0,0028 \\ \hline
1700       & 0,790& 0,007 & 4,88& 0,0029 \\ \hline
1800       & 0,925& 0,007 & 6,59& 0,0029 \\ \hline
 1900& 1,007& 0,007 & 7,74&0,0030\\\hline
\end{tabular}
\caption{Результаты измерений}
\end{table}
Проведем обработку результатов - определим по значениям яркостной температуры нити её термодинамическую температуру, используя график зависимости из справочных материалов, также посчитаем $W$. Результаты внесем в таблицу ниже:
\begin{table}[H]
\centering
\begin{tabular}{|c|c|c|c|c|c|c|}
\hline
$W$, Вт&  $\sigma_W$, мВт & $T$, $\text{С}^{\circ}$ & $\ln W$  & $\sigma_{\ln W}$ & $\ln T$  & $\sigma_{\ln T}$ \\ \hline
0.64& 2 & 908   & -0.443& 0,013      & 6,812 & 0,011      \\ \hline
0.90& 2 & 1017  & -0.101& 0,011      & 6,924 & 0,010      \\ \hline
1.09& 2 & 1125  & 0.087& 0,010      & 7,026 & 0,009      \\ \hline
1.37& 3 & 1233  & 0.318& 0,009      & 7,117 & 0,008      \\ \hline
1.76& 4 & 1342  & 0.564& 0,009      & 7,202 & 0,007      \\ \hline
2.33& 4 & 1450  & 0.845& 0,009      & 7,279 & 0,007      \\ \hline
2.71& 6 & 1558  & 0.995& 0,008      & 7,351 & 0,006      \\ \hline
3.72& 7 & 1667  & 1.313& 0,008      & 7,419 & 0,006      \\ \hline
3.86& 7 & 1775  & 1.349& 0,008      & 7,482 & 0,006      \\ \hline
6.10& 8 & 1883  & 1.808& 0,008      & 7,541 & 0,005      \\ \hline
\end{tabular}
\caption{Обработка результатов}
\end{table}
Построим график зависимости $\ln W (\ln T)$:

\begin{minipage}{1\textwidth}
    \centering
      \begin{tikzpicture}
        \begin{axis}[
          width=0.6\textwidth,
          height=,
          xlabel={$\ln{T}$},
          ylabel={$\ln{W}$},
          MyStyle1,
          ]
          \draw[red, domain=0:10] plot (\x,3.0532115353105045*\x-21.260771669670664);
          \addplot [only marks] plot [error bars/.cd,
                x dir=both, x explicit,
                y dir=both, y explicit,
                      ] table [
                      x=x, y=y,
                x error minus=errx, x error plus=errx, y error minus=erry, y error plus=erry, col sep=comma] {5.1.1.csv};
        \end{axis}
      \end{tikzpicture}
    \end{minipage}

% \begin{figure}[H]
%     \centering
%     \includegraphics[width=15cm]{pics/2.png}
% \end{figure}
Проведем наилучшую прямую по методу МНК, рассчитаем ее коэффициенты (все расчеты в этой работе проводятся на python с помощью соответствующих пакетов - numpy, pandas):
\begin{center}
    $k = (3.05 \pm 0.12)$, $b = (-21.27 \pm 0.86)$
\end{center}
По углу наклона графика можно оценить показатель степени температуры в законе Стефана-Больцмана. В нашем случае он получился примерно равным 3, при теоретическом значении 4.\\
Определим постоянную Стефана-Больцмана, используя значение термодинамической температуры 1883$^{\circ}$C и соответствующую мощность ($\varepsilon_T(1883) \approx 0,234$ и $S = 0,36$ см$^2$):
\begin{center}
    $\sigma_1 = \frac{W}{\varepsilon_T S T^4} = (3,89 \pm 0,06) \cdot 10^{-12}$ Вт/(см$^2 \cdot$ K$^4$)
\end{center}
Также можно определить постоянную Стефана-Больцмана, используя построенный график зависимости $\ln(W) = \ln(\varepsilon_T \sigma S) + n \ln(T)$:
\begin{center}
    $\ln(\varepsilon_t \sigma S) = (-21,3 \pm 1,2)$ \\
    $\sigma_2 = \frac{e^{-21,5}}{\varepsilon_T S} = (4,5 \pm 0,4) \cdot 10^{-12}$ Вт/(см$^2 \cdot$ K$^4$)
\end{center}
Табличное же значение постоянной Стефана-Больцмана равняется:
\begin{center}
    $\sigma_t = 5,67\cdot 10^{-12}$ Вт/(см$^2 \cdot$ K$^4$)
\end{center}
Как видно, оба значения достаточно близки к табличному.\\
Оценим значение постоянной Планка:
\begin{center}
    $h = \sqrt[3]{\frac{2 \pi^5 k_B^4}{15 c^2 \sigma}} = (7,14 \pm 0,42) \cdot 10^{-27}$ эрг$\cdot$с
\end{center}

\subsection*{Измерение яркостной температуры неоновой лампочки}
Термодинамическая температура неоновой лампочки примерно равна комнатной, и не соответствует её яркостной температуре ($\approx$ 900$^{\circ}$C). Дело в том, что неоновая лампочка в принципе не является моделью абсолютно чёрного или серого тела, и её излучение носит совершенно другую природу (переход электронов между энергетическими уровнями).
\section*{Выводы}
В ходе лабораторной работы удалось:
\begin{enumerate}
    % \item при помощи модели абсолютно чёрного тела (АТЧ) провести измерения температуры оптическим пирометром с исчезающей нитью и термопарой, при сравнении результатов обоих способов выяснилось, что измерения совпадают в пределах доверительных интервалов 
    % \item исследовать излучение накалённых тел с различной испускательной способностью: при проведении работы мы наблюдали, что для различных материалов с одинаковой термодинамической температурой их яркостная температура может не совпадать, что связано с различием коэффициента спектрального поглощения этих материалов
    \item определить постоянную Стефана-Больцмана двумя различными способами - через формулу $W = \varepsilon_T S \sigma T^4$ ($\sigma_1 = (3,89 \pm 0,06) \cdot 10^{-12}$ Вт/(см$^2 \cdot$ K$^4$)) и используя график зависимости $W(t)$ в логарифмическом масштабе ($\sigma_2 = (4,5 \pm 0,4) \cdot 10^{-12}$ Вт/(см$^2 \cdot$ K$^4$)), оба значения близки к табличному ($\sigma_t = 5,67\cdot 10^{-12}$ Вт/(см$^2 \cdot$ K$^4$)), но второй способ дает более верную оценку погрешности
    \item определить постоянную Планка $h = (7,14 \pm 0,42) \cdot 10^{-27}$ эрг$\cdot$с, полученной значение равно табличному ($h  = 6,626 \cdot 10^{-27}$ эрг$\cdot$с) в рамках доверительного интервала
    \item с помощью пирометра была определена "яркостная температура" неоновой лампочки, не совпадающая с термодинамической в силу того, что неоновая лампочка не является моделью АЧТ
\end{enumerate}
\end{document}